\documentclass[]{article}
\usepackage{lmodern}
\usepackage{amssymb,amsmath}
\usepackage{ifxetex,ifluatex}
\usepackage{fixltx2e} % provides \textsubscript
\ifnum 0\ifxetex 1\fi\ifluatex 1\fi=0 % if pdftex
  \usepackage[T1]{fontenc}
  \usepackage[utf8]{inputenc}
\else % if luatex or xelatex
  \ifxetex
    \usepackage{mathspec}
  \else
    \usepackage{fontspec}
  \fi
  \defaultfontfeatures{Ligatures=TeX,Scale=MatchLowercase}
\fi
% use upquote if available, for straight quotes in verbatim environments
\IfFileExists{upquote.sty}{\usepackage{upquote}}{}
% use microtype if available
\IfFileExists{microtype.sty}{%
\usepackage{microtype}
\UseMicrotypeSet[protrusion]{basicmath} % disable protrusion for tt fonts
}{}
\usepackage[margin=1in]{geometry}
\usepackage{hyperref}
\hypersetup{unicode=true,
            pdftitle={S4: The measurement and effect of an optimal offset for old interactions},
            pdfborder={0 0 0},
            breaklinks=true}
\urlstyle{same}  % don't use monospace font for urls
\usepackage{color}
\usepackage{fancyvrb}
\newcommand{\VerbBar}{|}
\newcommand{\VERB}{\Verb[commandchars=\\\{\}]}
\DefineVerbatimEnvironment{Highlighting}{Verbatim}{commandchars=\\\{\}}
% Add ',fontsize=\small' for more characters per line
\usepackage{framed}
\definecolor{shadecolor}{RGB}{248,248,248}
\newenvironment{Shaded}{\begin{snugshade}}{\end{snugshade}}
\newcommand{\AlertTok}[1]{\textcolor[rgb]{0.94,0.16,0.16}{#1}}
\newcommand{\AnnotationTok}[1]{\textcolor[rgb]{0.56,0.35,0.01}{\textbf{\textit{#1}}}}
\newcommand{\AttributeTok}[1]{\textcolor[rgb]{0.77,0.63,0.00}{#1}}
\newcommand{\BaseNTok}[1]{\textcolor[rgb]{0.00,0.00,0.81}{#1}}
\newcommand{\BuiltInTok}[1]{#1}
\newcommand{\CharTok}[1]{\textcolor[rgb]{0.31,0.60,0.02}{#1}}
\newcommand{\CommentTok}[1]{\textcolor[rgb]{0.56,0.35,0.01}{\textit{#1}}}
\newcommand{\CommentVarTok}[1]{\textcolor[rgb]{0.56,0.35,0.01}{\textbf{\textit{#1}}}}
\newcommand{\ConstantTok}[1]{\textcolor[rgb]{0.00,0.00,0.00}{#1}}
\newcommand{\ControlFlowTok}[1]{\textcolor[rgb]{0.13,0.29,0.53}{\textbf{#1}}}
\newcommand{\DataTypeTok}[1]{\textcolor[rgb]{0.13,0.29,0.53}{#1}}
\newcommand{\DecValTok}[1]{\textcolor[rgb]{0.00,0.00,0.81}{#1}}
\newcommand{\DocumentationTok}[1]{\textcolor[rgb]{0.56,0.35,0.01}{\textbf{\textit{#1}}}}
\newcommand{\ErrorTok}[1]{\textcolor[rgb]{0.64,0.00,0.00}{\textbf{#1}}}
\newcommand{\ExtensionTok}[1]{#1}
\newcommand{\FloatTok}[1]{\textcolor[rgb]{0.00,0.00,0.81}{#1}}
\newcommand{\FunctionTok}[1]{\textcolor[rgb]{0.00,0.00,0.00}{#1}}
\newcommand{\ImportTok}[1]{#1}
\newcommand{\InformationTok}[1]{\textcolor[rgb]{0.56,0.35,0.01}{\textbf{\textit{#1}}}}
\newcommand{\KeywordTok}[1]{\textcolor[rgb]{0.13,0.29,0.53}{\textbf{#1}}}
\newcommand{\NormalTok}[1]{#1}
\newcommand{\OperatorTok}[1]{\textcolor[rgb]{0.81,0.36,0.00}{\textbf{#1}}}
\newcommand{\OtherTok}[1]{\textcolor[rgb]{0.56,0.35,0.01}{#1}}
\newcommand{\PreprocessorTok}[1]{\textcolor[rgb]{0.56,0.35,0.01}{\textit{#1}}}
\newcommand{\RegionMarkerTok}[1]{#1}
\newcommand{\SpecialCharTok}[1]{\textcolor[rgb]{0.00,0.00,0.00}{#1}}
\newcommand{\SpecialStringTok}[1]{\textcolor[rgb]{0.31,0.60,0.02}{#1}}
\newcommand{\StringTok}[1]{\textcolor[rgb]{0.31,0.60,0.02}{#1}}
\newcommand{\VariableTok}[1]{\textcolor[rgb]{0.00,0.00,0.00}{#1}}
\newcommand{\VerbatimStringTok}[1]{\textcolor[rgb]{0.31,0.60,0.02}{#1}}
\newcommand{\WarningTok}[1]{\textcolor[rgb]{0.56,0.35,0.01}{\textbf{\textit{#1}}}}
\usepackage{graphicx,grffile}
\makeatletter
\def\maxwidth{\ifdim\Gin@nat@width>\linewidth\linewidth\else\Gin@nat@width\fi}
\def\maxheight{\ifdim\Gin@nat@height>\textheight\textheight\else\Gin@nat@height\fi}
\makeatother
% Scale images if necessary, so that they will not overflow the page
% margins by default, and it is still possible to overwrite the defaults
% using explicit options in \includegraphics[width, height, ...]{}
\setkeys{Gin}{width=\maxwidth,height=\maxheight,keepaspectratio}
\IfFileExists{parskip.sty}{%
\usepackage{parskip}
}{% else
\setlength{\parindent}{0pt}
\setlength{\parskip}{6pt plus 2pt minus 1pt}
}
\setlength{\emergencystretch}{3em}  % prevent overfull lines
\providecommand{\tightlist}{%
  \setlength{\itemsep}{0pt}\setlength{\parskip}{0pt}}
\setcounter{secnumdepth}{0}
% Redefines (sub)paragraphs to behave more like sections
\ifx\paragraph\undefined\else
\let\oldparagraph\paragraph
\renewcommand{\paragraph}[1]{\oldparagraph{#1}\mbox{}}
\fi
\ifx\subparagraph\undefined\else
\let\oldsubparagraph\subparagraph
\renewcommand{\subparagraph}[1]{\oldsubparagraph{#1}\mbox{}}
\fi

%%% Use protect on footnotes to avoid problems with footnotes in titles
\let\rmarkdownfootnote\footnote%
\def\footnote{\protect\rmarkdownfootnote}

%%% Change title format to be more compact
\usepackage{titling}

% Create subtitle command for use in maketitle
\newcommand{\subtitle}[1]{
  \posttitle{
    \begin{center}\large#1\end{center}
    }
}

\setlength{\droptitle}{-2em}

  \title{S4: The measurement and effect of an optimal offset for old interactions}
    \pretitle{\vspace{\droptitle}\centering\huge}
  \posttitle{\par}
    \author{true \\ true}
    \preauthor{\centering\large\emph}
  \postauthor{\par}
      \predate{\centering\large\emph}
  \postdate{\par}
    \date{2019-03-08}


\begin{document}
\maketitle

{
\setcounter{tocdepth}{3}
\tableofcontents
}
\newcommand{\R}{\mathbb{R}}
\newcommand{\E}{\mathbb{E}}
\newcommand{\V}{\mathbb{V}}
\newcommand{\Z}{\mathbb{Z}}
\newcommand{\N}{\mathbb{N}}
\newcommand{\B}{\mathscr{B}}
\newcommand{\W}{\mathscr{W}}
\newcommand{\ra}{\rightarrow}

\hypertarget{calculating-effect-size-for-old-interactions}{%
\subsubsection{Calculating effect size for old
interactions}\label{calculating-effect-size-for-old-interactions}}

For interactions taking place over extended periods of time, we must
take into account coevolutionary races will be put in check by some
other evolutionary forces. For parsimony, we assume this is performed by
some abstract ambient stabilizing selection acting separately on each
species. In particular, we extend the model derived in Appendix S1 so
that fitness of an individual with trait value \(z_i\) in species \(i\)
becomes

\[\mathscr{W}_i(z_i)=J_i\exp\left(-\frac{B_i}{2}(z_i-\bar z_j-\delta)^2\right)\exp\left(-\frac{A_i}{2}(z_i-\theta_i)^2\right)\]

for some \(J_i>0\). The equations of evolution, including our continuous
time approximation of drift, become

\[\begin{matrix}
d\bar z_1=G_1[B_1(\delta+\bar z_2-\bar z_1)+A_1(\theta_1-\bar z_1)]dt+\sqrt\frac{G_1}{N_1}dW_1\\
d\bar z_1=G_2[B_2(\delta+\bar z_1-\bar z_2)+A_2(\theta_2-\bar z_2)]dt+\sqrt\frac{G_2}{N_2}dW_2.
\end{matrix}\]

This process exhibits a bivariate normal stochastic equilibrium
distribution defined by

\[\mu_i=\frac{B_i(\mu_j+\delta)+A_i\theta_i}{B_i+A_i}\]

\[=\frac{B_iA_j+2B_iB_j}{(A_j+B_j)A_i+B_iA_j}\delta+\frac{(A_j+B_j)A_i\theta_i+B_iA_j\theta_j}{(A_j+B_j)A_i+B_iA_j}\]

\[\Sigma_{ii}=\frac{B_i\Sigma_{12}+\frac{1}{2N_i}}{B_i+A_i}\]

\[\Sigma_{12}=\frac{B_1(A_1+B_1)\frac{G_1}{2N_2}+
B_2(A_2+B_2)\frac{G_2}{2N_1}}{(A_1A_2+A_1B_2+A_2B_1)((A_1+B_1)G_1+(A_2+B_2)G_2)}.\]

Again, we find \(\Sigma_{ij}={\Sigma_{ij}}^*\). Hence, using the
Mahalanobis metric once more, we can calculate the effect of the optimal
offset on the traits of species that have coevolved in isolation for
long periods of time. Unfortunately, this measure is too cumbersome to
express even in an appendix. We therefore leave its full glory in the
associated \emph{Mathematica} notebook.

\hypertarget{measuring-the-optimal-offset-for-old-interactions-from-spatial-data}{%
\subsubsection{Measuring the optimal offset for old interactions from
spatial
data}\label{measuring-the-optimal-offset-for-old-interactions-from-spatial-data}}

To measure \(\delta\) indirectly for old interactions, we can employ the
maximum likelihood method of coevolution (MLC) introduced in Week \&
Nuismer (Week and Nuismer 2019). MLC requires estimates of local
population mean phenotypes across space and hence is not applicable for
single populations. However, with satisfactory data, MLC returns
reasonable estimates of biotic and abiotic selection strengths
(\(B_1,B_2,A_1,A_2\)) as well as the optimal offset (\(\delta\)). Thus,
MLC can easily be adapted to estimate the significance of the optimal
offset \(\delta\) using likelihood ratios analogous to those used to
estimate the significance of \(B_1\) and \(B_2\). In particular,
defining \(L_\delta\) the unconstrained likelihood of our model and
\(L_0\) the likelihood constrained to \(\delta=0\), we compute the
likelihood ratio statistic

\[\Lambda=2(\ln L_\delta-\ln L_0)\]

which is asymptotically \(\chi^2_1\) (with df\(=1\)). Writing
\(F_1(x), \ x>0\) the cumulative density function of the \(\chi^2_1\)
distribution, the \emph{p}-value associated with \(\Lambda\) is just
\(F_1(\Lambda)\). To perform this calculation, we adapt lines 169-248 of
the \emph{R} script \(\tt{estimate.R}\) found in the repository:

\(\tt{https://github.com/bobweek/measuring.coevolution}\)

This script depends also on \(\tt{functions.R}\) found in the same
repository. The adapted code is printed at the end of this document.

\hypertarget{r-code}{%
\subsubsection{\texorpdfstring{\emph{R} code}{R code}}\label{r-code}}

\begin{Shaded}
\begin{Highlighting}[]
\KeywordTok{require}\NormalTok{(mvtnorm)}
\KeywordTok{source}\NormalTok{(}\StringTok{"functions.R"}\NormalTok{)}

\CommentTok{#      #############      ########################################################}
\CommentTok{#      #           #}
\CommentTok{#      # Pauw data #      ########################################################}
\CommentTok{#      #           #}
\CommentTok{#      #############      ########################################################}

\NormalTok{snout <-}\StringTok{ }\KeywordTok{c}\NormalTok{(}\FloatTok{79.3}\NormalTok{,}\FloatTok{54.8}\NormalTok{,}\FloatTok{57.0}\NormalTok{,}\FloatTok{72.8}\NormalTok{,}\FloatTok{42.9}\NormalTok{,}\FloatTok{58.4}\NormalTok{,}\FloatTok{85.8}\NormalTok{,}\FloatTok{75.6}\NormalTok{)}
\NormalTok{tubes <-}\StringTok{ }\KeywordTok{c}\NormalTok{(}\FloatTok{77.0}\NormalTok{,}\KeywordTok{mean}\NormalTok{(}\KeywordTok{c}\NormalTok{(}\FloatTok{50.0}\NormalTok{,}\FloatTok{27.5}\NormalTok{)),}\FloatTok{49.1}\NormalTok{,}\FloatTok{52.6}\NormalTok{,}\FloatTok{41.1}\NormalTok{,}\FloatTok{58.6}\NormalTok{,}\FloatTok{60.8}\NormalTok{,}\FloatTok{68.0}\NormalTok{)}
\NormalTok{data <-}\StringTok{ }\KeywordTok{cbind}\NormalTok{(snout,tubes)}

\NormalTok{G1 <-}\StringTok{ }\KeywordTok{mean}\NormalTok{(}\KeywordTok{c}\NormalTok{(}\FloatTok{4.6}\NormalTok{,}\FloatTok{4.4}\NormalTok{,}\FloatTok{5.9}\NormalTok{,}\FloatTok{5.8}\NormalTok{,}\FloatTok{1.6}\NormalTok{,}\FloatTok{6.8}\NormalTok{,}\FloatTok{6.5}\NormalTok{,}\FloatTok{7.4}\NormalTok{)}\OperatorTok{^}\DecValTok{2}\NormalTok{)}
\NormalTok{G2 <-}\StringTok{ }\KeywordTok{mean}\NormalTok{(}\KeywordTok{c}\NormalTok{(}\FloatTok{5.9}\NormalTok{,}\FloatTok{4.9}\NormalTok{,}\FloatTok{4.3}\NormalTok{,}\FloatTok{4.7}\NormalTok{,}\FloatTok{2.5}\NormalTok{,}\FloatTok{7.5}\NormalTok{,}\FloatTok{4.6}\NormalTok{,}\FloatTok{5.5}\NormalTok{)}\OperatorTok{^}\DecValTok{2}\NormalTok{)}

\NormalTok{siblings <-}\StringTok{ }\KeywordTok{c}\NormalTok{(}\FloatTok{11.5}\NormalTok{,}\DecValTok{32}\NormalTok{,(}\DecValTok{38}\OperatorTok{+}\DecValTok{44}\NormalTok{)}\OperatorTok{/}\DecValTok{2}\NormalTok{)}

\NormalTok{ancepsAsolo <-}\StringTok{ }\KeywordTok{c}\NormalTok{(}\FloatTok{31.2}\NormalTok{,}\FloatTok{51.7}\NormalTok{)}
\NormalTok{Moegestorynchus <-}\StringTok{ }\KeywordTok{c}\NormalTok{(siblings,}\KeywordTok{mean}\NormalTok{(siblings)) }\CommentTok{# brevirostris, perplexus, braunsi}

\NormalTok{th1s <-}\StringTok{ }\NormalTok{Moegestorynchus}
\NormalTok{th2 <-}\StringTok{ }\KeywordTok{mean}\NormalTok{(ancepsAsolo)}
\NormalTok{n1 <-}\StringTok{ }\DecValTok{100}
\NormalTok{n2 <-}\StringTok{ }\DecValTok{1000}

\NormalTok{N <-}\StringTok{ }\KeywordTok{dim}\NormalTok{(data)[}\DecValTok{1}\NormalTok{]}
\NormalTok{m <-}\StringTok{ }\KeywordTok{colMeans}\NormalTok{(data)}
\NormalTok{S <-}\StringTok{ }\NormalTok{(N}\DecValTok{-1}\NormalTok{)}\OperatorTok{*}\KeywordTok{var}\NormalTok{(data)}\OperatorTok{/}\NormalTok{N}

\CommentTok{######}
\CommentTok{# ML #}
\CommentTok{######}

\NormalTok{B1s <-}\StringTok{ }\OtherTok{NULL}
\NormalTok{B2s <-}\StringTok{ }\OtherTok{NULL}
\NormalTok{A1s <-}\StringTok{ }\OtherTok{NULL}
\NormalTok{A2s <-}\StringTok{ }\OtherTok{NULL}
\NormalTok{ks <-}\StringTok{ }\OtherTok{NULL}
\NormalTok{ps <-}\StringTok{ }\OtherTok{NULL}
\ControlFlowTok{for}\NormalTok{(th1 }\ControlFlowTok{in}\NormalTok{ th1s)\{}
  
\NormalTok{  ML  <-}\StringTok{ }\KeywordTok{ml_sol}\NormalTok{(m[}\DecValTok{1}\NormalTok{],m[}\DecValTok{2}\NormalTok{],S[}\DecValTok{1}\NormalTok{,}\DecValTok{1}\NormalTok{],S[}\DecValTok{2}\NormalTok{,}\DecValTok{2}\NormalTok{],S[}\DecValTok{1}\NormalTok{,}\DecValTok{2}\NormalTok{],n1,n2,th1,th2,G1,G2)}
\NormalTok{  lnL <-}\StringTok{ }\KeywordTok{likelihood}\NormalTok{(data,m,S)}
  
\NormalTok{  A1s <-}\StringTok{ }\KeywordTok{c}\NormalTok{(A1s,ML}\OperatorTok{$}\NormalTok{A1)}
\NormalTok{  A2s <-}\StringTok{ }\KeywordTok{c}\NormalTok{(A2s,ML}\OperatorTok{$}\NormalTok{A2)}
\NormalTok{  B1s <-}\StringTok{ }\KeywordTok{c}\NormalTok{(B1s,ML}\OperatorTok{$}\NormalTok{B1)}
\NormalTok{  B2s <-}\StringTok{ }\KeywordTok{c}\NormalTok{(B2s,ML}\OperatorTok{$}\NormalTok{B2)}
\NormalTok{  ks  <-}\StringTok{ }\KeywordTok{c}\NormalTok{(ks,ML}\OperatorTok{$}\NormalTok{k)}
  
  \CommentTok{########}
  \CommentTok{# null #}
  \CommentTok{########}
  
\NormalTok{  ML_k0  <-}\StringTok{ }\KeywordTok{ml_sol_k0}\NormalTok{(m[}\DecValTok{1}\NormalTok{],m[}\DecValTok{2}\NormalTok{],S[}\DecValTok{1}\NormalTok{,}\DecValTok{1}\NormalTok{],S[}\DecValTok{2}\NormalTok{,}\DecValTok{2}\NormalTok{],S[}\DecValTok{1}\NormalTok{,}\DecValTok{2}\NormalTok{],n1,n2,th1,th2,G1,G2)}
  
\NormalTok{  A1_k0 <-}\StringTok{ }\NormalTok{ML_k0}\OperatorTok{$}\NormalTok{A1}
\NormalTok{  A2_k0 <-}\StringTok{ }\NormalTok{ML_k0}\OperatorTok{$}\NormalTok{A2}
\NormalTok{  B1_k0 <-}\StringTok{ }\NormalTok{ML_k0}\OperatorTok{$}\NormalTok{B1}
\NormalTok{  B2_k0 <-}\StringTok{ }\NormalTok{ML_k0}\OperatorTok{$}\NormalTok{B2}
  
\NormalTok{  mu_n <-}\StringTok{ }\KeywordTok{mu}\NormalTok{(A1_k0,A2_k0,B1_k0,B2_k0,}\DecValTok{0}\NormalTok{,n1,n2,th1,th2)}
\NormalTok{  S_n <-}\StringTok{ }\KeywordTok{SIGMA}\NormalTok{(A1_k0,A2_k0,B1_k0,B2_k0,}\DecValTok{0}\NormalTok{,n1,n2,th1,th2,G1,G2)}
\NormalTok{  lnL0  <-}\StringTok{ }\KeywordTok{likelihood}\NormalTok{(data,mu_n,S_n)}
  
  \CommentTok{# chisqr}
\NormalTok{  ps <-}\StringTok{ }\KeywordTok{c}\NormalTok{(ps,}\DecValTok{1}\OperatorTok{-}\KeywordTok{pchisq}\NormalTok{(}\DecValTok{2}\OperatorTok{*}\NormalTok{(lnL}\OperatorTok{-}\NormalTok{lnL0),}\DecValTok{1}\NormalTok{))}
  
\NormalTok{\}}

\NormalTok{pval <-}\StringTok{ }\KeywordTok{max}\NormalTok{(ps)}
\end{Highlighting}
\end{Shaded}

\hypertarget{refs}{}
\leavevmode\hypertarget{ref-Week2019}{}%
Week, Bob, and Scott L. Nuismer. 2019. ``The Measurement of Coevolution
in the Wild.'' Edited by Frederick Adler. \emph{Ecology Letters},
February. Wiley. \url{https://doi.org/10.1111/ele.13231}.


\end{document}
